\documentclass[10pt,twocolumn,letterpaper]{article}
\usepackage[T1]{fontenc}
\usepackage[utf8]{inputenc}
\usepackage{amsmath,amssymb,amsthm,mathtools,bm,braket}
\usepackage{mathptmx}
\renewcommand{\ttdefault}{pcr}
\usepackage[margin=0.68in,columnsep=0.26in,headheight=14pt,headsep=0.15in,footskip=0.28in]{geometry}
\usepackage{graphicx,booktabs,array,tabularx,longtable,multirow}
\usepackage{microtype,fancyhdr,caption,titlesec,enumitem,float,placeins}
\usepackage[most]{tcolorbox}
\usepackage[hidelinks,bookmarksnumbered=true,hypertexnames=false]{hyperref}
\usepackage{xurl}
\setlength{\parindent}{1em}
\setlength{\parskip}{0.16em}
\setlength{\emergencystretch}{2em}
\setlength{\abovedisplayskip}{6pt plus2pt minus2pt}
\setlength{\belowdisplayskip}{6pt plus2pt minus2pt}
\captionsetup{font=small,labelfont=bf,justification=raggedright,singlelinecheck=false}
\titleformat{\section}{\large\bfseries}{\thesection.}{0.45em}{}
\titleformat{\subsection}{\normalsize\bfseries}{\thesubsection}{0.45em}{}
\titleformat{\subsubsection}{\normalsize\bfseries}{\thesubsubsection}{0.45em}{}
\titlespacing*{\section}{0pt}{11pt plus2pt minus2pt}{5pt}
\titlespacing*{\subsection}{0pt}{8pt plus2pt minus2pt}{3pt}
\pagestyle{fancy}\fancyhf{}
\fancyhead[L]{\footnotesize SAM Volume I | Chapter \ChapterRoman}
\fancyhead[R]{\footnotesize Brady | Edition of 25 October 2026}
\fancyfoot[C]{\footnotesize\thepage}
\renewcommand{\headrulewidth}{0.3pt}
\newtheorem{theorem}{Theorem}[section]
\newtheorem{proposition}{Proposition}
\newtheorem{lemma}[theorem]{Lemma}
\newtheorem{corollary}[theorem]{Corollary}
\theoremstyle{definition}
\newtheorem{definition}[theorem]{Definition}
\newtheorem{assumption}[theorem]{Assumption}
\newtheorem{example}[theorem]{Example}
\theoremstyle{remark}
\newtheorem{remark}[theorem]{Remark}
\newcommand{\dd}{\mathrm{d}}
\newcommand{\ii}{\mathrm{i}}
\newcommand{\Tr}{\operatorname{Tr}}
\newcommand{\Span}{\operatorname{span}}
\allowdisplaybreaks[2]
\raggedbottom
\pdfmapfile{+cm.map}
\pdfmapfile{+cmextra.map}
\pdfmapfile{+symbols.map}
\pdfmapfile{+psnfss.map}
\pdfmapfile{+utm.map}
\pdfmapfile{=ps2pk35.map}
\pdfmapfile{+rsfs.map}

\newcommand{\ChapterRoman}{I}
\hypersetup{pdftitle={SAM Volume I: Substrate and accumulation, the concepts},pdfauthor={Sean Brady}}
\begin{document}
\fancyhead[R]{\footnotesize Brady | Draft of 30 October 2026}
\twocolumn[{
\begin{center}
{\small SAM VOLUME I\par}
\vspace{0.4em}
{\LARGE\bfseries Substrate and accumulation: the concepts\par}
\vspace{0.4em}
{\large Chapter I\par}
\vspace{0.5em}
{\large Sean Brady\par}
{\small Originator and conceptual director\par}
{\small Research development with OpenAI ChatGPT and Codex; concept chapter drafted with Claude Code\par}
{\small Draft of 30 October 2026 for the first edition\par}
\end{center}
\begin{minipage}{0.97\textwidth}\noindent\textbf{Abstract.} This chapter is the conceptual introduction to the volume. It states the single sentence of the Substrate Accumulation Model, the three objects that sentence names (the floor $A_0$, the lift $A(r)$ and the ceiling $1/\pi$), and the two readings of the lift that a clock and a ruler make. It states the agreement theorem for the static exterior, the two numbers SAM keeps between any pair of endpoints where most frameworks keep one, the SAM law for the line of sight, and the clock in motion. It ends where SAM parts from general relativity: at $A=1$, a closed sphere with no interior, which every traveler approaches and none crosses. Every statement is labeled as derived, declared or open, and the closing section records the verification behind each.\end{minipage}
\vspace{0.9em}}]


\section{One sentence and a picture}
\label{sec:c1-sentence}

The Substrate Accumulation Model (SAM) rests on one sentence, in the words of its originator:
\begin{quote}
\textbf{``Matter displaces space, and that displacement accumulates.''}
\end{quote}
Everything in this volume is an attempt to give that sentence a value, a law of composition and a set of consequences that can be checked. The sentence has three parts, and each becomes an object in Section~\ref{sec:c1-three}. \emph{Space} is the thing displaced, and SAM calls it the substrate and assigns it a value. \emph{Matter displaces} it, and the displacement of one source at a distance $r$ is the lift $A(r)$. \emph{The displacement accumulates}: along any line through space the displacements of all sources add, and the sum along the longest lines has a ceiling.

The picture that preceded the sentence is the aquarium, a visualization from 2023 that stands as it was envisioned. It is not a model to be tested at every angle and no equation is written from it. It is the way to show a student the idea.
\begin{quote}
An indestructible aquarium the size of a room, filled completely with water, with no room to add a single atom and no possibility for a single atom to escape. A small spherical rock the size of a softball grows into existence in the middle. With the water having nowhere to go, a halo of ``denser'' water forms around the rock. A small pebble floating in a straight line towards the rock enters the halo, slows, and its path bends around the rock; once it leaves the halo it continues in a straight path as before.
\end{quote}
The pebble floats in something. When the value of that something changes, because matter has displaced it, the pebble behaves differently. The bending is the gradient, the pull of the rock. The slowing is what a distant observer sees of anything passing through accumulated substrate, and Section~\ref{sec:c1-home} returns to it. Nothing in the picture concerns clocks, and the water does not move.

\section{The floor}
\label{sec:c1-floor}

SAM's first claim about space is that it is made of something and has a value. The value where nothing has displaced it is the floor, written $A_0$. In the originator's words:
\begin{quote}
The idea of SAM is that space (the substrate) is made of something, and SAM attempts to apply a value to it. $A_0$ is theorized as the minimum value for space itself to exist. If matter displaces space and that displacement accumulates, then, like land above sea level, it lifts from a floor, and that floor is $A_0$.
\end{quote}
The number is fixed in Chapters~II and~III by the counting convention of the finite ledger,
\begin{equation}
A_0=\frac{1}{12\pi}=0.026525\ldots,
\label{eq:c1-floor}
\end{equation}
and three properties follow from its being a floor rather than a field.

First, a constant has no gradient. A value greater than zero everywhere, however large, defines no direction and no rest frame; nothing can be measured against it. Only differences in $A$ act. Everything local in this volume, every clock rate, every deflection, every delay, depends on how far above the floor a point sits relative to another, never on the floor's own height. This is why $A_0$ drops out of every static-exterior observable in Section~\ref{sec:c1-agreement} and why GPS, which has nanosecond sensitivity, sees no trace of it.

Second, the floor is large compared with what the Earth produces. The Earth's lift at its surface is about $1.4\times10^{-9}$; the floor is more than ten million times that. The land is high and the hills on it are low. The sea-level picture is exact in this sense: a hill's height above the surrounding land is what a walker feels, not the land's height above the sea.

Third, the floor stands alone only in the vastest and emptiest regions. The Earth, the Sun and everything in the Milky Way produce accumulation, so no local measurement is ever made at the floor. Where the floor does act is where differences cannot: in the count of how much space there is, which is the cosmology of Chapters~VIII and~XII, and in the ceiling of Section~\ref{sec:c1-three}. The convention adopted throughout this volume is the source-only reading, in which the lift $A(r)$ is measured from the floor and the floor is not added to it inside any clock or ruler; the alternative, background-normalized reading rescales the lift by $1/(1-A_0)$ and is labeled as a convention wherever it appears.

\section{The three objects}
\label{sec:c1-three}

\paragraph{The lift.} A spherical source of mass $M$ displaces space around it. The displacement at areal radius $r$ is
\begin{equation}
A(r)=\frac{r_s}{r},\qquad r_s=\frac{2GM}{c^2},
\label{eq:c1-lift}
\end{equation}
largest against the surface of the source and falling with distance. Chapter~III derives this profile from the spherical accumulation kernel; it is the response of the substrate to a source and is not assumed. At $r=r_s$ the lift reaches unity, which is the top of the hill and the subject of Section~\ref{sec:c1-home}. The lift of several sources adds at a point: accumulation at the Earth's surface is $A_0$ plus the Earth's own lift, plus the Sun's at one astronomical unit, plus the Galaxy's, each small, each present.

\paragraph{The road.} Accumulation is also what a line of sight passes through. For a route $\Gamma$ between two points, define
\begin{equation}
\Sigma[\Gamma]=\int_\Gamma a\,\dd\ell,
\label{eq:c1-road}
\end{equation}
the accumulation looked through, where $a$ is the local lift and $\dd\ell$ the measurable length element along the route. Chapter~IV calls this the road operator. It is a property of the path, not of its endpoints, and Section~\ref{sec:c1-two} shows that this one distinction carries most of SAM's content.

\paragraph{The ceiling.} Along the longest roads, those that cross the expanding universe, the accumulation looked through does not grow without limit. Chapter~XII writes the line-of-sight accumulation to redshift $z$ as
\begin{equation}
A_{\rm los}(z)=\frac{1}{\pi}\left[1-(1+z)^{-3}\right],
\label{eq:c1-ceiling}
\end{equation}
which rises from zero and saturates at $1/\pi$. The ceiling is the end of the longest road. It is the counting convention's number, $1/\pi=12A_0 R$ with $R$ the count, and it takes no part in anything static. It is named here because the three objects belong together: the floor is where the substrate starts, the lift is what one source does to it, the ceiling is how much of it a line of sight can ever look through.

\section{Two readings of the lift}
\label{sec:c1-readings}

A lift is a number at a point. Physics enters when an instrument reads it, and SAM declares two readings.

\paragraph{The clock.} A clock held at rest where the lift is $a$ runs, relative to a clock at the floor, at the rate
\begin{equation}
N(a)=\sqrt{1-a}.
\label{eq:c1-clock}
\end{equation}
Chapter~IV builds this from the source-to-clock interface. Two clocks at different lifts compare as $N(a_2)/N(a_1)$, a function of the two endpoint values alone.

\paragraph{The ruler.} A ruler laid radially where the lift is $a$ reads, per unit of areal radius,
\begin{equation}
\dd\ell=\frac{\dd r}{\sqrt{1-a}},
\label{eq:c1-ruler}
\end{equation}
while a ruler laid along a sphere reads the areal length $r\,\dd\theta$. Chapter~IV states the radial reading as a declaration tested by the Cassini conjunction, not as a derivation, and Appendix~A keeps the derivation open. This is the one place in the static exterior where SAM declares rather than derives, and the volume says so each time it is used.

The two readings together are a metric,
\begin{equation}
\dd s^2=-(1-a)\,c^2\dd t^2+\frac{\dd r^2}{1-a}+r^2\dd\Omega^2,
\label{eq:c1-metric}
\end{equation}
and this is the static exterior of a spherical source in SAM.

\section{The agreement theorem and what it does and does not say}
\label{sec:c1-agreement}

Equation~\eqref{eq:c1-metric} is the Schwarzschild exterior. Every static observable of a spherical vacuum therefore takes, in SAM, the value general relativity assigns it, exactly and not to first order: the gravitational redshift between two heights, the deflection of light $4GM/(c^2 b)$ at impact parameter $b$, the Shapiro delay, the perihelion advance, the geodetic precession, the photon sphere at $a=2/3$ and the shadow of diameter $3\sqrt3\,r_s$, the innermost stable circular orbit at $a=1/3$. In the parametrized post-Newtonian language, $\gamma=\beta=1$. The Cassini measurement of $\gamma-1=(2.1\pm2.3)\times10^{-5}$ is a test of the ruler declaration~\eqref{eq:c1-ruler} and it is passed.

What the theorem says is that SAM is not refuted by any static measurement in the solar system or at the photon sphere of a black hole. What it does not say is that SAM is general relativity. General relativity reproduces Newton's numbers in the weak field and nobody calls general relativity Newton's work. Equal numbers in a regime do not make equal theories; the theories differ where their accountings differ, and the chapters that follow are about where SAM's accounting differs. There are three such places, and this chapter introduces each: the two numbers between endpoints (Section~\ref{sec:c1-two}), the road taken to cosmological length (Section~\ref{sec:c1-law}), and the boundary at $A=1$ (Section~\ref{sec:c1-home}). The reader should hold to one rule throughout: SAM's results are SAM's results. Where they coincide with general relativity the coincidence is stated; SAM is never described as ``just GR'' in one place and as departing from it in another by convenience.

\section{Two numbers between two endpoints}
\label{sec:c1-two}

\subsection{The magic bell}

The question that separated the clock from the road was practical: how can accumulation be reconciled with the need to adjust the clocks in GPS satellites? The thought experiment that answered it is the magic bell, in the originator's words:
\begin{quote}
On one large planet near a supermassive black hole sits a mirror clock with one observer. On a second, tiny planet in a void sits only an observer with a magic bell that rings exactly when light strikes the mirror. The question is: does the observer hear the bell and see the light strike the mirror in unison?
\end{quote}
The question is not answered in this volume, by decision. Its value is the two things it forces apart. The mirror clock sits deep in accumulation and the bell's observer sits in the void, so their clocks run at different rates; that is the clock, a comparison of two endpoints. And the light that carries the strike from one planet to the other has to cross the accumulated substrate in between; that is the road, a property of the path. The two are not the same quantity, and the result of the thought experiment, in the originator's words, is a statement about the second in terms of the first:
\begin{quote}
If $A$ is not equal at two endpoints, accumulation must exist between them.
\end{quote}

\subsection{The theorem and its failed converse}

In the volume's terms the result is a theorem, and the proof is one line: $a$ is continuous along any route, so if $a_1\ne a_2$ then $a>0$ on a set of positive length, and $\Sigma=\int a\,\dd\ell>0$. Between two endpoints SAM therefore keeps two numbers where most frameworks keep one. The first is the measurable distance of the route,
\begin{equation}
L[\Gamma]=\int_\Gamma \dd\ell,
\end{equation}
what a ruler laid along it reads. The second is the accumulation looked through, $\Sigma[\Gamma]$ of Eq.~\eqref{eq:c1-road}. The clock reads the endpoints; the road reads the path.

The converse fails, and the failure is as important as the theorem. Two endpoints can have the same lift with a great deal of accumulation between them. The Cassini conjunction of 2002 is exactly that case: the Earth at one astronomical unit and the spacecraft at 8.4 had lifts that differed by less than one part in $10^{8}$, so the clock comparison between them was a few parts in $10^9$ and unchanged by the conjunction, while the route between them swung within $1.6$ solar radii of the Sun and the accumulation looked through rose to $\Sigma=13.3\,r_s$, a one-way delay of $131\,\mu$s. The clock did not move; the road did. That is the Shapiro delay read as SAM reads it, and it is the sentence ``accumulation is looked through, not only sat in'' made measurable.

\subsection{GPS and Shapiro, the two instances}

The same displacement gives GPS its two corrections, and they are applied separately because they are the two numbers. The clock correction compares the satellite's endpoint with the ground's: at $26\,560$~km the gravitational term is $+45.7\,\mu$s per day, the velocity term (Section~\ref{sec:c1-motion}) is $-7.2$, and the net is $+38.5$, or $+38.6$ once the ground clock's own rotation is counted, which is the figure engineered into the satellites. The road correction is the accumulation looked through along the signal's path, $\Sigma/c$, and it is $42$~ps for a zenith route and $62$~ps for a horizon route, small because the Earth's hill is low, present because the route crosses it. One displacement, two readings, two corrections. The Shapiro delay is the road correction grown to planetary scale; GPS is the clock correction made daily. Between them they are the first two chapters of any account of accumulation, and the volume gives them a chapter of their own (Chapter~IV) beyond this introduction.

\section{The SAM law}
\label{sec:c1-law}

Take the road to its full length. In the originator's words:
\begin{quote}
\textbf{``The farther you look out into space, the more accumulation you are looking through.''}
\end{quote}
This is the SAM law, and it is the sentence of Section~\ref{sec:c1-sentence} read along a line of sight. It is not a challenge to Hubble's law. The redshift of distant galaxies grows with their distance, and that is an observation; the SAM law offers a different explanation of the same observation, that a photon from far away has crossed more accumulated substrate before its distance is inferred. Chapter~XII writes the explanation as the road law~\eqref{eq:c1-ceiling}, with the ceiling $1/\pi$ at the end of the longest road, and Chapter~VIII carries the inventory that follows from it. The law makes a prediction that Hubble's law alone does not: the clock number and the road number of a distant source are different quantities, so the luminosity distance and the angular-diameter distance need not be related as a pure expansion relates them, and Chapter~XII names the differential test that would decide it. Here the law is introduced as the third object's law: the ceiling is where the SAM law ends.

\section{The clock in motion}
\label{sec:c1-motion}

A clock has two reasons to run slow: where it sits and how fast it moves. Chapter~IV's exact law for a clock at lift $a$ moving with speed $v$ relative to the frame in which the field is static is
\begin{equation}
\frac{\dd\tau}{\dd t}=\sqrt{1-a}\,\sqrt{1-\frac{v^2}{c^2}},
\label{eq:c1-motion}
\end{equation}
which to first order is $1-a/2-\beta^2/2$ with $\beta=v/c$: position and speed add. For a circular orbit the two terms are locked together, $\beta^2=a/2$, and the rate is $\sqrt{1-3a/2}$ exactly, with the family of circular orbits ending at the photon sphere $a=2/3$. The GPS numbers of Section~\ref{sec:c1-two} are this one line evaluated once: the orbit's $a$ gives $+45.7\,\mu$s a day, its $\beta^2$ gives $-7.2$. Along a slightly eccentric orbit the same line gives the periodic correction every receiver applies, with the constant the GPS specification prints reproduced to eleven digits. The speed in Eq.~\eqref{eq:c1-motion} is relative to the frame in which the field is static. It is not a speed through the floor, because, as Section~\ref{sec:c1-floor} said, the floor is a constant and a constant defines nothing to move against; a clock knows the gradient it sits in and its motion relative to what makes that gradient, and nothing else. The clock in motion is the third reading of the lift, after the clock at rest and the ruler, and Chapter~IV develops it.

\section{The horizon: where SAM parts from general relativity}
\label{sec:c1-home}

\subsection{The boundary}

The lift reaches unity at $r=r_s$. In SAM, accumulation beyond completion, $A>1$, is not a state that can be written: there is no register past full. A black hole is therefore a closed sphere with no interior. There is no event on the far side of $A=1$ for a crossing to become, and so nothing crosses it. This is a theoretical departure, and it is the headline of the volume, not something to be hidden. It is also structural. It is not a statement about speed or time, and no engine, however fast, changes it. Every observer approaches the boundary; none crosses. Chapter~V develops the approach, and Chapters~XII and~XIII what a completed sphere is and does.

\subsection{The approach, in SAM's accounting}

How long the approach takes is the second question. The outside observer's account is shared by every framework: the coordinate time to reach $r_s$ diverges, and the traveler fades down the redshift rather than winking out. The traveler's own account is where SAM's accounting differs, and SAM's accounting is the one this volume carries.

The distance that decides whether a crossing completes is the measurable road, the distance a ruler built from the field reads, published by the originator in July 2026 and carried into Chapter~V. Along the approach that road is
\begin{equation}
L(r)=\int_r^{r_0}\frac{\dd\rho}{1-A(\rho)},
\label{eq:c1-measurable-road}
\end{equation}
which is $c$ times the time light takes and which diverges logarithmically as $r\downarrow r_s$: every nine added to $A$ adds $r_s\ln 10$ of road, forever, while the metric's own length to the horizon stays finite (from $2r_s$ it is $2.30\,r_s$ in all). An astronaut cannot travel at $c$. The least time his own clock can spend on each nine of road is therefore the road at light speed, $r_s\ln10/c$, and a traveler of any mass spends more. The total is unbounded in either case. He experiences time normally and writes in his journal every day, but the accumulated road in front of him grows without bound, and he never arrives at $A=1$. In the originator's published words: ``The astronaut does not run out of clock. He runs out of road, there is always more measurable distance ahead, and no $A>1$ state behind it to write into. He arrives at $A=1$. He never crosses it.''

What SAM has not yet derived is the law that fixes how fast a traveler of a given mass covers the measurable road, and so how much longer than the light-speed floor each nine takes him. That law sets the pace of the approach. It does not bear on whether he arrives, which is settled once the measurable road is the physical distance and nothing with mass moves along it at $c$.

\subsection{The other accounting, and what the two share}

General relativity's accounting is the free-fall geodesic of the metric~\eqref{eq:c1-metric}, with the metric's proper length taken as the physical distance. In that accounting a traveler released from rest far away and passing $2r_s$ reaches $r_s$ after $1.219\,r_s/c$ of his own time, a few microseconds for a stellar-mass hole and under a day for the largest known, and crosses without any local mark of the passage. SAM does not accept that accounting as the physical one, for the reason Section~\ref{sec:c1-two} gave: the metric's length and the measurable road are two different numbers between the same endpoints, and at the boundary they separate completely.

The two accountings agree far from the boundary and split at it, and the split is the clock against the road. Time dilation is relational and is not experienced; in both accountings the traveler's clock ticks normally and dilation exists only in comparison with another clock. Road accumulation is local and is experienced; in SAM's accounting there is always more measurable distance ahead, and that is a fact of the traveler's own frame. The same split shows in the traveler's speed as a nearby static observer would measure it: in the free-fall accounting it climbs toward $c$ as $A\to1$; in SAM's accounting it falls toward zero as the road thickens under him, a tenth of $c$ at two nines, a thousandth at six. That is the pebble of Section~\ref{sec:c1-sentence} slowing in the halo. Table~\ref{tab:c1-approach} sets the two accountings side by side for three masses.

\begin{table}[t]
\centering\small
\caption{The approach to $A=1$ from $2r_s$ in the two accountings. GR: the free-falling traveler's own time to the horizon, in total. SAM: the least possible own time per nine added to $A$, the road at light speed; a traveler of any mass spends more on each nine, and the total is unbounded.}
\label{tab:c1-approach}
\begin{tabular}{@{}lccc@{}}
\toprule
 & $r_s/c$ & GR: to $r_s$ & SAM: per nine, at least \\
\midrule
one solar mass & $9.85\,\mu$s & $12\,\mu$s & $23\,\mu$s \\
Sagittarius A$^*$ & $42$~s & $52$~s & $98$~s \\
M87$^*$ & $17.8$~h & $22$~h & $41$~h \\
\bottomrule
\end{tabular}
\end{table}

\subsection{What is derived, what is declared, what is open}

The structural claim, that nothing crosses $A=1$, follows from the ledger: no $A>1$ state exists to be written into. The divergence of the measurable road, Eq.~\eqref{eq:c1-measurable-road}, is a theorem. That the measurable road is the physical distance is SAM's premise, published and carried here, and with it the conclusion that no traveler with mass arrives at $A=1$ in his own time. What is open is the law of motion for mass on the measurable road: it fixes the pace of the approach for a given mass, and Chapter~V names it as the volume's obligation, together with the requirement that it leave every weak-field result, from the muon's lifetime to the GPS clock, exactly as measured. Nothing outside can decide between the accountings: the clock, the shadow and every weak-field test are blind to the difference. The traveler can, and in SAM he can tell by the road, never by his clock.

\section{How the volume proceeds}
\label{sec:c1-map}

Chapters~II and~III build the substrate from a finite ledger of source histories: the retention rule, the accumulation count, the floor $A_0$ as the count's minimum, and the spherical kernel that turns a source into the lift $A(r)$. Chapter~IV reads the lift with clocks and signals: the clock law, the ruler declaration and its Cassini test, the road operator, the clock in motion, and GPS as the daily instance. Chapter~V takes the lift to unity: the three questions of the approach, the criterion for a finite or infinite own time, the road-speed postulate and its cost, and the boundary at $A=1$. Chapter~VI carries the quantum routes by which records are read. Part~2 extends the program: horizon thermodynamics (VII), the cosmic inventory and the Hubble cascade under the road law (VIII), the thermal history (IX), carrier radiation (XI), the cosmological fate and the completed boundary (XII), and the completed sphere and its chain (XIII). Appendix~A keeps the register of what is derived, what is declared and what remains to be assigned.

\section{Status and verification}
\label{sec:c1-status}

This chapter is conceptual by design and its claims are of three kinds. \emph{Derived}: the agreement theorem and its classical numbers; the two-numbers theorem and its failed converse; the clock in motion, its circular-orbit form and the GPS corrections; the divergence of the measurable road; the criterion for the traveler's own time. \emph{Declared}: the radial ruler reading, tested by Cassini; the source-only convention for the floor; the road-speed postulate at the boundary. \emph{Open}: the derivation of the ruler from the source action; the law of motion for mass on the measurable road; the differential test of the road law against distance data. The question of the magic bell is open by decision.

Each derived or declared statement above was run in the verification form, with predictions sealed before computation, in the Courtroom~II record: the static exterior theorem and its classical numbers (\texttt{voli\_metric/01}); the two numbers, the Cassini conjunction and the GPS road term (\texttt{voli\_clock\_road/01}); the clock in motion, GPS from one line, the closure line and the eccentricity correction (\texttt{voli\_clock\_motion/01} and \texttt{/04}); the approach to $A=1$ under the declared readings and through the measurable road (\texttt{g11555r/01} and \texttt{/02}). These are reproductions with standard tools, not independent verification in the sense of Appendix~A, and nothing in this chapter is classified.

Credit for the ideas is recorded as it fell. Sean Brady: the sentence, the aquarium, the floor as the minimum value of space, the magic bell and its result, the SAM law, the measurable road, the boundary at $A=1$ as a closed sphere with no interior, and that no traveler with mass arrives there. OpenAI ChatGPT and Codex: the mathematical development of Chapters~II to~VI from which the derived statements are drawn. Claude (web): the proposal of the agreement theorem as a chapter statement and the one-line clock in motion. Claude Code: the corrections to those proposals, the computations in the verification record, and the drafting of this chapter under direction. Directed by Sean Brady.

\end{document}
